\documentclass[12pt]{article}
\input{preamble}

\pagestyle{fancy}
\fancyhf{}

\rhead{Nørre Gymnasium\\
2.e
}


%Husk at rette modul og dato!
\lhead{Matematik A\\
Modul 23
}
\chead{1/4/2021
}

\cfoot{Side \thepage \hspace{1pt} af \pageref{LastPage}}

\begin{document}

%Udfyld afsnit herunder og lav til egen Latex-fil

%Kopier følgende til overskrift:

%\begin{center}
%\Huge
%Aflevering 1
%\end{center}
%\section*{Opgave 1}
%\stepcounter{section}

\begin{center}
\Huge
Omdregningslegemer
\end{center}

\section*{Omdregningslegemer}
\stepcounter{section}

Har vi en funktion $f$, der på et interval $[a,b]$ opfylder, at $f(x)>0$, og roterer den omkring $x$-aksen i rummet, så afgrænser denne rotation et \textit{omdregningslegeme}. Vi kan se dette på Fig. \ref{fig:omdregning}
\begin{figure}[H]
\centering
\begin{tikzpicture}
\begin{axis}[axis lines = middle, xmin = -2.5, xmax = 5.5, ymin = -20, ymax = 20,
ticks = none
]
\addplot[samples = 1000, color = blue!40, domain = -1:4] {x^3 - 4*x^2 + x + 10};
\draw[dashed] (axis cs:-1,0) -- (axis cs:-1,4);
\draw[dashed] (axis cs:4,0) -- (axis cs:4,14);
\addplot[samples = 1000, color = blue!40, domain = -1:4] {-(x^3 - 4*x^2 + x + 10)};
\draw[dashed, color = blue!40] (axis cs:-1,0) ellipse (0.25cm and 0.7cm);
\end{axis}
\end{tikzpicture}
\end{figure}

\section*{Opgave 1}
Løs følgende integraler med integration ved substitution:
\begin{align*}
&1)   \  \int e^{2x+1}x \intd x  &&2) \ \int \frac{e^{\sqrt{x}}}{\sqrt{x}} \intd x   \\
&3)   \   \int \frac{6x^2}{2\sqrt{x^3+1}} \intd x &&4) \  \int -(x^3-\frac{1}{2}x)\sin(x^4-x^2) \intd x   \\
\end{align*}
\section*{Opgave 2}
\begin{enumerate}[label=\roman*)]
\item Bestem arealet mellem funktionen $f(x) = e^{x^3}\cdot x^2$ og $x$-aksen på intervallet $[0,2]$.
\item Bestem arealet mellem funktionerne $f(x) = e^{\sqrt{x}}/\sqrt{x}$ og $g(x) = 2x+3$ på intervallet $[4,16]$
\end{enumerate}
\end{document}
